% Modern editorial, markup and figure/layout contributions: CC-BY-SA-4.0.
% Historical text and illustrations: Leonhard Euler (1744), public domain.
% Component scope and upstream attribution: see RIGHTS.md.
% Original opening page is misprinted 309; logical pagination is 311.
\subsection*{Additamentum II.}
\textit{De motu projectorum in medio non resistente, per Methodum maximorum ac minimorum determinando.}
\originalsection{1}\sourcepage{311}\footnote{The source page is printed 309; the original contents and the page sequence identify it as 311.}Quoniam omnes naturæ effectus sequuntur quandam maximi minimive legem; dubium est nullum, quin in lineis curvis, quas corpora projecta, si a viribus quibuscunque sollicitentur, describunt, quæpiam maximi minimive proprietas locum habeat. Quænam autem sit ista proprietas, ex principiis metaphysicis a priori definire non tam facile videtur: cum autem has ipsas curvas, ope Methodi directæ, determinare liceat; hinc, debita adhibita attentione, id ipsum, quod in istis curvis est maximum vel minimum, concludi poterit. Spectari autem potissimum debet effectus a viribus sollicitantibus oriundus; qui cum in motu corporis genito consistat, veritati consentaneum videtur hunc ipsum motum, seu potius aggregatum omnium motuum qui in corpore projecto insunt, minimum esse debere. Quæ conclusio etsi non satis confirmata videatur, tamen, si eam cum veritate jam a priori nota consentire ostendero, tantum consequetur pondus, ut omnia dubia quæ circa eam suboriri queant penitus evanescant. Quin etiam cum ejus veritas fuerit evicta, facilius erit in intimas Naturæ leges atque causas finales inquirere; hocque assertum firmissimis rationibus corroborare.

\originalsection{2}Sit massa corporis projecti $=M$, ejusque, dum spatiolum $=ds$ emetitur, celeritas debita altitudini $=v$; erit quantitas motus corporis in hoc loco $=M\sqrt v$; quæ per ipsum spatiolum $ds$ multiplicata, dabit $Mds\sqrt v$ motum corporis collectivum per spatiolum $ds$. Jam dico lineam a corpore descriptam \sourcepage{312}ita fore comparatam, ut, inter omnes alias lineas iisdem terminis contentas, sit $\int Mds\sqrt v$, seu, ob $M$ constans, $\int ds\sqrt v$ minimum. Quod si autem curva quæsita tanquam esset data spectetur, ex viribus sollicitantibus celeritas $\sqrt v$ per quantitates ad curvam pertinentes definiri, ideoque ipsa curva per Methodum maximorum ac minimorum determinari potest. Cæterum hæc expressio ex quantitate motus petita æque ad vires vivas traduci poterit; posito enim tempusculo, quo elementum $ds$ percurritur, $=dt$; quia est $ds=dt\sqrt v$, fiet $\int ds\sqrt v=\int vdt$; ita ut, in curva a corpore projecto descripta, summa omnium virium vivarum, quæ singulis temporis momentis corpori insunt, sit minima. Quamobrem neque ii qui vires per ipsas celeritates, neque illi qui per celeritatum quadrata æstimari oportere statuunt, hic quicquam quo offendantur reperient.

\originalsection{3}Primum igitur, si corpus a nullis prorsus viribus sollicitari ponamus, ejus quoque celeritas, ad quam hic solum attendo (directionem enim ipsa Methodus maximorum \& minimorum complectetur), nullam patietur alterationem; eritque ideo $v$ quantitas constans, puta $=b$. Hinc corpus a nullis viribus sollicitatum, si utcunque projiciatur, ejusmodi describet lineam, in qua sit $\int ds\sqrt b$ vel $\int ds=s$ minimum. Via ergo hæc, inter omnes iisdem terminis contentas, ipsa erit minima; atque adeo recta: prorsus uti prima Mechanicæ principia postulant. Hunc quidem casum non adeo hic affero, quo principium meum confirmari putem; quamcunque enim, loco celeritatis $\sqrt v$, aliam assumsissem functionem ipsius $v$, eadem prodiisset via recta; verum a casibus simplicissimis incipiendo facilius ipsa consensus ratio intelligi poterit.

\originalsection{4}Progredior ergo ad casum gravitatis uniformis, seu quo corpus projectum ubique, secundum directiones ad horizontem normales, deorsum sollicitetur a vi constante acceleratrice $=g$. Sit $AM$ curva, quam corpus in hac hypothesi describit, sumatur recta verticalis $AP$ pro axe, ac ponatur abscissa $AP=x$, applicata $PM=y$, \& elementum curvæ $Mm=ds$; erit ergo, ex natura sollicitationis, $dv=gdx$, \& $v=a+gx$. (Fig.~26.) Hinc \sourcepage{313}curva ita erit comparata, ut in ea sit $\int ds\sqrt{(a+gx)}$ minimum. Ponatur $dy=pdx$, ut sit $ds=dx\sqrt{(1+pp)}$, atque minimum esse debet $\int dx\sqrt{(a+gx)(1+pp)}$; quæ expressio cum forma generali $\int Zdx$ comparata dat $Z=\sqrt{(a+gx)(1+pp)}$; quare, cum positum sit $dZ=Mdx+Ndy+Pdp$, erit $N=0$ \& $P=\frac{p\sqrt{(a+gx)}}{\sqrt{(1+pp)}}$. Quia ergo valor differentialis est $N-\frac{dP}{dx}$; ob $N=0$, fiet præsenti casu $dP=0$, \& $P=\sqrt C$. Habebitur ergo
\[\sqrt C=\frac{p\sqrt{(a+gx)}}{\sqrt{(1+pp)}}=\frac{dy\sqrt{(a+gx)}}{ds};\]
unde fit $Cdx^2+Cdy^2=dy^2(a+gx)$, \& $dy=\frac{dx\sqrt C}{\sqrt{(a-C+gx)}}$; quæ integrata dat $y=\frac2g\sqrt{C(a-C+gx)}$.

\originalsection{5}Manifestum quidem est hanc æquationem esse pro Parabola. At ejus consensum cum veritate attentius considerasse juvabit. Primum ergo patet tangentem hujus curvæ esse horizontalem, seu $dx=0$; ubi est $a-C+gx=0$. Cum igitur principium abscissarum $A$ ab arbitrio nostro pendeat, sumatur id in hoc ipso loco, fietque $C=a$; tum vero ipse axis per hoc punctum curvæ summum transeat, ita ut, posito $x=0$, fiat simul $y=0$. His consideratis, æquatio pro curva erit hæc $y=2\sqrt{\frac{ax}{g}}$; quam non solum patet esse pro Parabola; sed etiam, cum celeritas in puncto $A$ sit $\sqrt a$, altitudo $CA$, ex qua corpus labendo ab eadem vi $g$ sollicitatum eam ipsam acquirit celeritatem, qua in puncto $A$ movetur, erit $=\frac ag$; hoc est, quartæ parametri parti æquatur; prorsus uti ex doctrina motus projectorum per Methodum directam intelligitur.

\originalsection{6}Sollicitetur, ut ante, corpus ubique verticaliter deorsum, at ipsa vis sollicitans non sit constans, sed pendeat utcunque ab altitudine $CP$. Scilicet posita abscissa $CP=x$, sit vis qua corpus in $M$ deorsum nititur $=X$ functioni cuicunque ipsius $x$. Si ergo vocetur applicata $PM=y$, elementum arcus $Mm$
\sourcepage{314}$=ds$, \& $dy=pdx$; erit $dv=Xdx$, \& $v=A+\int Xdx$; unde minimum esse debet hæc expressio $\int dx\sqrt{(A+\int Xdx)(1+pp)}$, ex qua pro curva descripta $AM$ obtinebitur hæc æquatio
\[\sqrt C=\frac{p\sqrt{(A+\int Xdx)}}{\sqrt{(1+pp)}}\quad\&\quad p=\frac{\sqrt C}{\sqrt{(A-C+\int Xdx)}}=\frac{dy}{dx};\]
seu $y=\int\frac{dx\sqrt C}{\sqrt{(A-C+\int Xdx)}}$. Tangens ergo curvæ erit horizontalis ubi $\int Xdx=C-A$. Hæc vero eadem æquatio trajectoriæ corporis per Methodum directam reperitur.

\originalsection{7}Sollicitetur nunc corpus in $M$ a duabus viribus, altera horizontali $=Y$ secundum directionem $MP$, altera verticali $=X$ secundum directionem $MQ$. Sit autem $X$ functio quæcunque rectæ verticalis $MQ=CP=x$, \& $Y$ functio quæcunque applicatæ $PM=y$. (Fig.~27.) Positis ergo ut ante $dy=pdx$, erit $dv=-Xdx-Ydy$, fietque $v=A-\int Xdx-\int Ydy$; unde minimum esse debet hæc formula $\int dx\sqrt{(1+pp)(A-\int Xdx-\int Ydy)}$. Differentietur $\sqrt{(1+pp)(A-\int Xdx-\int Ydy)}$, atque prodibit
\[\begin{aligned}
&-\frac{Xdx\sqrt{(1+pp)}}{2\sqrt{(A-\int Xdx-\int Ydy)}}-\frac{Ydy\sqrt{(1+pp)}}{2\sqrt{(A-\int Xdx-\int Ydy)}}\\
&+\frac{pdp\sqrt{(A-\int Xdx-\int Ydy)}}{\sqrt{(1+pp)}}.
\end{aligned}\]
Hinc posito
\[N=\frac{-Y\sqrt{(1+pp)}}{2\sqrt{(A-\int Xdx-\int Ydy)}}\quad\&\quad P=\frac{p\sqrt{(A-\int Xdx-\int Ydy)}}{\sqrt{(1+pp)}};\]
erit pro curva quæsita hæc æquatio $0=N-\frac{dP}{dx}$, seu $Ndx=dP$. Hinc ergo fit
\[\begin{aligned}
\frac{-Ydx\sqrt{(1+pp)}}{2\sqrt{(A-\int Xdx-\int Ydy)}}
&=\frac{dp\sqrt{(A-\int Xdx-\int Ydy)}}{(1+pp)\sqrt{(1+pp)}}\\
&\quad-\frac{pXdx-pYdy}{2\sqrt{(1+pp)(A-\int Xdx-\int Ydy)}},
\end{aligned}\]
seu
\[\frac{dp\sqrt{(A-\int Xdx-\int Ydy)}}{(1+pp)\sqrt{(1+pp)}}=\frac{Xdy-Ydx}{2\sqrt{(1+pp)(A-\int Xdx-\int Ydy)}};\]
ideoque $\frac{2dp}{1+pp}=\frac{Xdy-Ydx}{A-\int Xdx-\int Ydy}$. Hanc æquationem veritati esse consentaneam patebit, si loco $A-\int Xdx-\int Ydy$ \sourcepage{315}ponatur $v$, erit enim
\[\frac{2vdp}{(1+pp)^{3:2}}=\frac{Xdy-Ydx}{\sqrt{(1+pp)}}.\]
At est radius osculi $r=-\frac{(1+pp)^{3:2}dx}{dp}$, quo introducto est $\frac{2v}{r}=\frac{Ydx-Xdy}{ds}$; ubi est $\frac{2v}{r}$ vis corporis centrifuga, \& $\frac{Ydx-Xdy}{ds}$ exprimit vim normalem ex viribus sollicitantibus ortam; quarum virium æqualitas utique in omni motu projectorum locum habet.

\originalsection{8}Æquatio autem inventa $\frac{2dp}{1+pp}=\frac{Xdy-Ydx}{A-\int Xdx-\int Ydy}$ ita generaliter est integrabilis, si multiplicetur per $\frac{p(A-\int Xdx-\int Ydy)}{1+pp}$; fiet enim
\[\frac{2pdp(A-\int Xdx-\int Ydy)}{(1+pp)^2}-\frac{ppXdx+Ydy}{1+pp}=0;\]
quæ integrata dat
\[\frac{-p^2\int Xdx+\int Ydy-A}{1+pp}=C,\]
seu $\int Ydy-p^2\int Xdx=A+C+Cpp$, unde $p=\frac{\sqrt{(B+\int Ydy)}}{\sqrt{(C+\int Xdx)}}$; posito $B$ pro $-A-C$. Cum ergo sit $p=\frac{dy}{dx}$, erit
\[\int\frac{dy}{\sqrt{(B+\int Ydy)}}=\int\frac{dx}{\sqrt{(C+\int Xdx)}},\]
æquatio pro curva quæsita, in qua variabiles $x$ \& $y$ sunt a se invicem separatæ. Vel si constantes $B$ \& $C$ in negativas convertantur, erit
\[\int\frac{dy}{\sqrt{(B-\int Ydy)}}=\int\frac{dx}{\sqrt{(A-\int Xdx)}}.\]
Ex quibus etsi curvæ constructio facilis habetur, tamen æquationes algebraicæ, quoties quidem in ipsis continentur, non tam facile eruuntur. Sint $X$ \& $Y$ functiones similes \& quidem potestates ipsarum $x$ \& $y$, ita ut sit
\[\int\frac{dy}{\sqrt{(b^n-y^n)}}=\int\frac{dx}{\sqrt{(a^n-x^n)}},\]
quæ æquatio, si $n=1$, præbet Parabolam; sin $n=2$, Ellipsin centrum in $C$ habentem: etsi hoc casu utraque integratio quadraturam Circuli \sourcepage{316}requirit. Verisimile ergo videtur etiam aliis casibus, quibus neutra integratio succedit, curvas algebraicas satisfacere; quarum autem inveniendarum Methodus adhuc desideratur.

\originalsection{9}Urgeatur corpus $M$ perpetuo versus punctum fixum secundum directionem $MC$, vi quæ sit ut functio quæcunque distantiæ $MC$. Positis ut ante $CP=x$, $PM=y$, \& $dy=pdx$; sit $CM=\sqrt{(x^2+y^2)}=t$, atque sit $T$ ea functio ipsius $t$, quæ exprimit vim centripetam. Resolvatur hæc vis in laterales secundum $MQ$ \& $MP$, erit vis trahens secundum $MQ=\frac{Tx}{t}$; \& vis secundum $MP=\frac{Ty}{t}$; ex quibus oritur acceleratio $dv=-\frac{Txdx}{t}-\frac{Tydy}{t}=-Tdt$, ob $xdx+ydy=tdt$; unde fit $v=A-\int Tdt$. Quamobrem minimum esse debet hæc expressio $\int dx\sqrt{(1+pp)(A-\int Tdt)}$. Jam, secundum Regulæ præceptum, differentietur quantitas $\sqrt{(1+pp)(A-\int Tdt)}$, prodibitque
\[-\frac{Tdv\sqrt{(1+pp)}}{2\sqrt{(A-\int Tdt)}}+\frac{pdp\sqrt{(A-\int Tdt)}}{\sqrt{(1+pp)}}.\]
Ob $dt=\frac{xdx+ydy}{t}$, erit ergo
\[N=\frac{-Ty\sqrt{1+pp}}{2t\sqrt{(A-\int Tdt)}}\quad\&\quad P=\frac{p\sqrt{(A-\int Tdt)}}{\sqrt{(1+pp)}};\]
ex quibus efficitur æquatio pro curva $Ndx=dP$, quæ præbet,
\[-\frac{Tydx\sqrt{(1+pp)}}{2t\sqrt{(A-\int Tdt)}}=\frac{dp\sqrt{(A-\int Tdt)}}{(1+pp)\sqrt{(1+pp)}}-\frac{pTdt}{2\sqrt{(1+pp)(A-\int Tdt)}},\]
hæcque reducta abibit in istam,
\[\frac{T(xdy-ydx)}{2t(A-\int Tdt)}=\frac{dp}{1+pp}.\]

\originalsection{10}Quamvis hæc æquatio quatuor contineat litteras diversas, tamen debita dexteritate integrari potest. Cum enim sit $ydy+xdx=tdt=pydx+xdx$, erit $dx=\frac{tdt}{x+py}$ \& $dy=\frac{ptdt}{x+py}$; qui valores in æquatione substituti dabunt
\[\frac{(px-y)Tdt}{2(x+py)(A-\int Tdt)}=\frac{dp}{1+pp},\quad\text{seu}\quad\frac{Tdt}{2(A-\int Tdt)}=\frac{dp(py+x)}{(1+pp)(px-y)}.\]
\sourcepage{317}Harum expressionum utraque per logarithmos est integrabilis, est enim $\int\frac{Tdt}{2(A-\int Tdt)}=-\frac12l(A-\int Tdt)$, \& $\int\frac{dp(x+py)}{(1+pp)(px-y)}$ resolvitur in $\int\frac{xdp}{px-y}-\int\frac{pdp}{1+pp}=l\frac{px-y}{\sqrt{(1+pp)}}$; ita ut sit
\[\frac C{\sqrt{(A-\int Tdt)}}=\frac{px-y}{\sqrt{(1+pp)}};\]
qua æquatione declaratur, celeritatem corporis in $M$, quæ est $=\sqrt{(A-\int Tdt)}$, esse reciproce ut perpendiculum ex $C$ in tangentem demissum; quæ est proprietas palmaria horum motuum.

\originalsection{11}Hoc vero idem Problema commodius resolvi potest ipsam rectam $CM$ pro altera variabili assumendo. Verum Methodus supra tradita non postulat, ut ambæ variabiles sint coordinatæ orthogonales, dummodo sint ejusmodi binæ quantitates quibus determinatis simul curvæ punctum determinetur. Hanc ob causam, non liceret distantiam $CM$ cum perpendiculo ex $C$ in tangentem demisso pro binis illis variabilibus accipere; quoniam etiamsi detur \& distantia a centro \& perpendiculum in tangentem, hinc tamen locus puncti curvæ non definitur. Nihil autem impedit, quo minus distantia $CM$, \& arcus circuli $BP$ centro $C$ descripti, in locum duarum variabilium substituantur; quia dato arcu $BP$, \& distantia $CM$ curvæ punctum $M$ æque determinatur ac per coordinatas orthogonales. (Fig.~28.) Hac ergo annotatione usus Methodi multo latius extenditur, quam alioquin videri queat.

\originalsection{12}Sit igitur distantia corporis a centro $MC=x$, \& vis qua corpus ad centrum $C$ sollicitatur sit $=X$ functioni cuicunque ipsius $x$. Centro $C$, radio pro lubitu assumpto $BC=c$, describatur circulus, cujus arcus $BP$ teneat locum alterius variabilis $y$, ita ut sit $Pp=dy=pdx$. Ex vi autem sollicitante est $dv=-Xdx$, unde $v=A-\int Xdx$. Centro $C$, radio $CM=x$, describatur arculus $Mn$, erit $mn=dx$; \& $CP:Pp=CM:Mn$, unde fit $Mn=\frac{pxdx}{c}$, \& elementum spatii $Mm=dx\sqrt{(1+\frac{ppxx}{cc})}$. Quamobrem minimum esse \sourcepage{318}debet hæc formula $\int dx\sqrt{(A-\int Xdx)(1+\frac{ppxx}{cc})}$, ex qua oritur valor differentialis
\[\frac1{dx}d.\frac{pxx\sqrt{(A-\int Xdx)}}{c\sqrt{(cc+ppxx)}},\]
qui, per Regulam, nihilo æqualis positus, præbebit hanc æquationem:
\[\sqrt C=\frac{pxx\sqrt{(A-\int Xdx)}}{c\sqrt{(cc+ppxx)}},\]
seu $Cc^4+Cccppxx=(A-\int Xdx)ppx^4$, ex qua fit
\[p=\frac{cc\sqrt C}{\sqrt{((A-\int Xdx)x^4-Cccxx)}}=\frac{cc\sqrt C}{x\sqrt{((A-\int Xdx)xx-Ccc)}}\]
seu
\[dy=\frac{ccdx\sqrt C}{x\sqrt{((A-\int Xdx)xx-Ccc)}},\]
quæ eadem æquatio etiam per Methodum directam invenitur.

\originalsection{13}Ex his igitur casibus perfectissimus consensus principii hic stabiliti cum veritate elucet: utrum autem iste consensus in casibus magis complicatis locum quoque sit habiturus, dubium superesse potest. Quamobrem quam late pateat istud principium diligentius erit investigandum, quo plus ipsi non tribuatur quam ejus natura permittit. Ad hoc explicandum, omnis motus projectorum in duo genera distribui debet; quorum altero celeritas corporis, quam in quovis loco habet, a solo situ pendet; ita ut, si ad eundem situm revertatur, eandem quoque sit recuperaturum celeritatem; quod evenit, si corpus vel ad unum vel ad plura centra fixa trahatur viribus, quæ sint ut functiones quæcunque distantiarum ab his centris. Ad alterum genus refero eos projectorum motus, quibus celeritas corporis per solum locum in quo hæret non determinatur; id quod usu venit, vel si centra illa ad quæ corpus sollicitatur fuerint mobilia, vel si motus fiat in medio resistente. Hac facta divisione; notandum est, quoties motus corporis ad prius genus pertineat, hoc est, si corpus non solum ad unum sed ad quotcunque centra fixa sollicitetur viribus quibuscunque, toties in motu hoc summam omnium motuum elementarium fore minimam.

\originalsection{14}Hoc ipsum autem postulat indoles Propositionis: dum enim, inter datos terminos, ea quæritur curva, in qua sit $\int ds\sqrt v$ minimum; eo ipso assumitur, celeritatem corporis in utroque \sourcepage{319}termino eandem esse, quæcunque curva corporis viam constituat. Quotcunque autem fuerint centra virium fixa, celeritas corporis in quovis loco $M$, exprimitur functione determinata ambarum variabilium $CP=x$, \& $PM=y$. Sit igitur $v$ functio quæcunque ipsarum $x$ \& $y$, ita ut sit $dv=Tdx+Vdy$; atque videamus, an principium nostrum veram exhibiturum sit projectoriam corporis. (Fig.~27.) Cum autem sit $dv=Tdx+Vdy$; corpus perinde movebitur, ac si sollicitetur in $M$ a duabus viribus, altera $T$ in directione abscissis $x$ parallela, altera vero $V$ in directione parallela applicatis $y$, ex quibus oritur vis tangentialis $=\frac{Tdx+Vdy}{ds}$, \& vis normalis $=\frac{-Vdx+Tdy}{ds}$. Debet autem, ex natura motus liberi, esse
\[\frac{2v}{r}=\frac{-Vdx+Tdy}{ds}=\frac{-V+Tp}{\sqrt{(1+pp)}};\]
ad quam æquationem si Methodus maximorum ac minimorum deducat, erit utique principium nostrum veritati conforme.

\originalsection{15}Cum igitur, per hoc principium, debeat esse $\int dx\sqrt{v(1+pp)}$ minimum, differentietur quantitas $\sqrt{v(1+pp)}$, atque, ob $dv=Tdx+Vdy$, orietur:
\[\frac{Tdx\sqrt{(1+pp)}}{2\sqrt v}+\frac{Vdy\sqrt{(1+pp)}}{2\sqrt v}+\frac{pdp\sqrt v}{\sqrt{(1+pp)}},\]
ex quo obtinetur pro curva quæsita sequens æquatio, secundum præcepta tradita,
\[\frac{Vdx\sqrt{(1+pp)}}{2\sqrt v}=d.\frac{p\sqrt v}{\sqrt{(1+pp)}}=\frac{dp\sqrt v}{(1+pp)^{1:2}}+\frac{p(Tdx+Vdy)}{2\sqrt{v(1+pp)}};\]
seu
\[\frac{-dp\sqrt v}{(1+pp)^{3:2}}=\frac{Tpdx-Vdx}{2\sqrt{v(1+pp)}}.\]
At est radius osculi in $M=\frac{-(1+pp)dx\sqrt{(1+pp)}}{dp}$; qui si ponatur $=r$, erit $\frac{2v}{r}=\frac{Tp-V}{\sqrt{(1+pp)}}$; omnino uti per Methodum directam invenitur. Dummodo ergo vires sollicitantes ita fuerint comparatæ, ut eæ reduci queant ad duas vires $T$ \& $V$, secundum directiones coordinatis $x$ \& $y$ parallelas sollicitantes, quæ sint ut functiones quæcunque harum variabilium $x$ \& $y$, tum semper in \sourcepage{320}curva descripta erit motus corporis per omnia elementa collectus minimus.

\originalsection{16}Tam late ergo hoc principium patet, ut solus motus a resistentia medii perturbatus excipiendus videatur; cujus quidem exceptionis ratio facile perspicitur, propterea quod hoc casu corpus per varias vias ad eundem locum perveniens non eandem acquirit celeritatem. Quamobrem, sublata omni resistentia in motu corporum projectorum, perpetuo hæc constans proprietas locum habebit, ut summa omnium motuum elementarium sit minima. Neque vero hæc proprietas in motu unius corporis tantum cernetur, sed etiam in motu plurium corporum conjunctim; quæ quomodocunque in se invicem agant, tamen semper summa omnium motuum est minima. Quod, cum hujusmodi motus difficulter ad calculum revocentur, facilius ex primis principiis intelligitur, quam ex consensu calculi secundum utramque Methodum instituti. Quoniam enim corpora, ob inertiam, omni status mutationi reluctantur; viribus sollicitantibus tam parum obtemperabunt, quam fieri potest, siquidem sint libera: ex quo efficitur, ut, in motu genito, effectus a viribus ortus minor esse debeat, quam si ullo alio modo corpus vel corpora fuissent promota. Cujus ratiocinii vis, etiamsi nondum satis perspiciatur; tamen, quia cum veritate congruit, non dubito quin, ope principiorum sanioris Metaphysicæ, ad majorem evidentiam evehi queat; quod negotium aliis, qui Metaphysicam profitentur, relinquo.
